\chapter*{Technical Abstract}
\addcontentsline{toc}{chapter}{Technical Abstract}

This project investigates whether three-dimensional active acoustic localisation can be performed using a structured neuromorphic model inspired by bat echolocation. The task is to estimate target distance, azimuth, and elevation from the echo produced by a known emitted call. This is relevant to edge sensing because localisation systems often face constraints on computation, latency, memory, and power. Low-power operation is not measured directly in this project, so it is treated as a motivation rather than a claimed result. The technical question addressed here is whether biological structure provides a useful way to design an interpretable spiking neural network (SNN) sensing system.

The model is built as a complete simulated echolocation pipeline. A frequency-modulated call is emitted into a controlled acoustic environment, and the returning binaural echo is generated with distance-dependent delay, propagation attenuation, interaural timing and level differences, and an elevation-dependent comb-filter spectral cue. The received waveform is processed by an IIR cochlear filterbank and a dynamic leaky integrate-and-fire spike encoder, producing a frequency-organised event-based representation. This representation is then passed through three parallel pathway models. Distance is estimated using echo delay, corollary discharge, onset extraction, and delay coincidence. Azimuth is represented using both interaural time-difference coincidence and interaural level-difference opponent comparison. Elevation is represented using a spectral-template model that compares the observed spectrum with expected comb-filter transfer functions. These pathways are simplified functional abstractions rather than detailed anatomical simulations, but each is tied to a specific acoustic cue and biological computation.

The pathway outputs are represented as one-dimensional neural populations. Direct centre-of-mass readouts, calibrated readouts, and superior-colliculus-style continuous-attractor readouts are compared as population decoders. The continuous attractor neural network (CANN) readouts are interpreted conservatively: they provide stabilising coordinate readouts, but they do not act as universal optimisers. Their effect is pathway-dependent, so the final model preserves both the raw pathway populations and the fixed readout estimates.

The main trainable component is a small SNN fusion stage. It receives raw distance, azimuth, and elevation populations, fixed coordinate readouts, and confidence features such as population activation and spike-count statistics. Two trainable variants are evaluated: a direct readout from the cached pathway features and a residual readout that learns a correction to the fixed pathway estimate. Additional SNN baselines are trained from less structured inputs, including raw binaural waveforms and cochlear spike rasters. These baselines test whether the hand-designed auditory pathways provide useful feature extraction beyond generic SNN processing of the original signal.

Performance is evaluated in the constrained operating region of 0.25--5~m and $\pm45^\circ$ azimuth/elevation. Errors are reported as distance mean absolute error, angular mean absolute error, Euclidean position error, and a dimensionless combined error formed from normalised distance and angular errors. Reported ranges are bootstrap 95\% confidence intervals over a 400-sample held-out test set.

In noise-free conditions, the fixed pathway model already performs functional localisation. The fixed CANN/calibrated readout obtains a combined error of 0.046 [0.041, 0.051] and Euclidean mean absolute error of 0.244 [0.204, 0.287]~m. The residual pathway-feature SNN gives the best noise-free result, reducing combined error to 0.022 [0.020, 0.024] and Euclidean error to 0.125 [0.110, 0.140]~m, with distance, azimuth, and elevation errors of 0.014 [0.013, 0.015]~m, 2.62 [2.36, 2.89]$^\circ$, and 0.27 [0.25, 0.29]$^\circ$ respectively.

Environmental noise exposes the main limitation of the fixed model. When a 50~dB environmental noise field is added before binaural and elevation filtering, the fixed CANN/calibrated baseline suffers a large elevation failure: its combined error rises to 0.263 [0.246, 0.279] and Euclidean error to 1.577 [1.447, 1.707]~m. The residual SNN substantially corrects this failure, reducing combined error to 0.055 [0.051, 0.059] and Euclidean error to 0.297 [0.266, 0.332]~m. With both environmental noise and delayed echo components, the residual SNN remains stable, with combined error 0.055 [0.050, 0.060] and Euclidean error 0.310 [0.277, 0.348]~m. In these noisy conditions, feature ablation indicates that the residual SNN depends strongly on raw elevation and distance population features, while CANN scalar readouts contribute less when ablated in isolation.

The input-only SNN baselines support the same interpretation. Raw waveform SNNs perform poorly despite receiving the original acoustic signal. Cochlear-raster SNNs perform substantially better, showing that early auditory preprocessing is useful. However, the best overall performance is obtained when the SNN receives the full set of biologically motivated pathway features. The project therefore demonstrates a hybrid design principle: fixed biomimetic pathways provide interpretable and useful feature extraction, while a small trainable SNN provides contextual integration and residual correction.

The current system remains a constrained single-object, single-pulse prototype. It uses simplified acoustic cues, functional rather than biophysical neural models, and hand-tuned pathway parameters. Nevertheless, the model establishes a complete foundation for extending neuromorphic echolocation toward continuous multi-pulse tracking, active feedback control, more realistic acoustic scenes, partially trainable pathways, and direct hardware or energy evaluation.
